\item Remítase al problema 16 y la figura P14.16. Se construirá un hidrómetro con una barra cilíndrica flotante. Se colocarán nueve marcas a lo largo de la barra para indicar densidades que tengan valores de $0.98 \text{ g/cm}^3$, $1.00 \text{ g/cm}^3$, $1.02 \text{ g/cm}^3$, $1.04 \text{ g/cm}^3$, \dots, $1.14 \text{ g/cm}^3$. La hilera de marcas comenzará $0.200 \text{ cm}$ desde el extremo superior de la barra y terminará $1.80 \text{ cm}$ desde el extremo superior. (a) ¿Cuál es la longitud requerida de la barra? (b) ¿Cuál debe ser su densidad promedio? (c) ¿Las marcas deberían estar igualmente espaciadas? Explique su respuesta.
